\section{AMC 12 12B 2009}
        \begin{enumerate}
        	\item (\textit{AMC 12 12B 2009}) Each morning of her five-day workweek, Jane bought either a $50$-cent muffin or a $75$-cent bagel. Her total cost for the week was a whole number of dollars. How many bagels did she buy?


 $\textbf{(A) } 1\qquad\textbf{(B) } 2\qquad\textbf{(C) } 3\qquad\textbf{(D) } 4\qquad\textbf{(E) } 5$

 \textbf{Answer: }B



\textbf{Solution 1 (Observations: Replacements)}

If Jane bought one more bagel but one fewer muffin, then her total cost for the week would increase by $25$ cents.

 If Jane bought $1$ bagel, then she bought $4$ muffins. Her total cost for the week would be $75\cdot1+50\cdot4=275$ cents, or $2.75$ dollars. Clearly, she bought one more bagel but one fewer muffin at a total cost of $3.00$ dollars. Therefore, she bought $\boxed{\textbf{(B) } 2}$ bagels.

 



\textbf{Solution 2 (Observations: Answer Choices)}


            \begin{itemize}
        \item If Jane bought $1$ bagel, then she bought $4$ muffins. Her total cost for the week would be $75\cdot1+50\cdot4=275$ cents, or $2.75$ dollars. So, $\textbf{(A)}$ is incorrect.
            \end{itemize}
            
            \begin{itemize}
        \item If Jane bought $2$ bagels, then she bought $3$ muffins. Her total cost for the week would be $75\cdot2+50\cdot3=300$ cents, or $3.00$ dollars. So, $\boxed{\textbf{(B) } 2}$ is correct.
            \end{itemize}
            For completeness, we will check $\textbf{(C)},\textbf{(D)},$ and $\textbf{(E)}$ too. If you decide to use this solution on the real test, then you will not need to do that, as you want to save more time.

 
            \begin{itemize}
        \item If Jane bought $3$ bagels, then she bought $2$ muffins. Her total cost for the week would be $75\cdot3+50\cdot2=325$ cents, or $3.25$ dollars. So, $\textbf{(C)}$ is incorrect.
            \end{itemize}
            
            \begin{itemize}
        \item If Jane bought $4$ bagels, then she bought $1$ muffin. Her total cost for the week would be $75\cdot4+50\cdot1=350$ cents, or $3.50$ dollars. So, $\textbf{(D)}$ is incorrect.
            \end{itemize}
            
            \begin{itemize}
        \item If Jane bought $5$ bagels, then she bought $0$ muffins. Her total cost for the week would be $75\cdot5+50\cdot0=375$ cents, or $3.75$ dollars. So, $\textbf{(E)}$ is incorrect.
            \end{itemize}
            



\textbf{Solution 3 (Arithmetic)}

In this solution, all amounts are in the unit of cents.

 Note that the amount spent on muffins must end in either $00$ or $50,$ and the amount spent on bagels must end in one of $00,25,50,$ or $75.$ Furthermore, the number of muffins bought and the number of bagels bought must sum to $5.$

 We have two possible cases:

 
            \begin{enumerate} 
        \item The amounts spent on muffins and bagels both end in $00.$The number of muffins bought must be even, and the number of bagels bought must be a multiple of $4.$ 

 In this case, there are no solutions.  

 
        \item The amounts spent on muffins and bagels both end in $50.$The number of muffins bought must be odd, and the number of bagels bought must be $2$ more than a multiple of $4.$ 

 In this case, the only solution is $3$ muffins and $2$ bagels, from which the answer is $\boxed{\textbf{(B) } 2}.$

            \end{enumerate}
            



\textbf{Solution 4 (Casework, similar to others)}

We require
 \[\frac{1}{2} m + \frac{3}{4} b \in \mathbb{Z}\].We approach with standard casework. Obviously, letting $m=0,b=5$ or vice versa does not work. If $m = 1,b=4$ the right side is integral whereas the left side is not on the LHS. The same applies for the versa of that case. If $m = 2. b = 3$ then the left side is integral but the right is not of the LHS. However, when $m=3, b=2$, we have that $\frac{3}{2} + \frac{3}{2}  = 3$ is indeed integral. The answer is $\boxed{2}$.

 


	\item (\textit{AMC 12 12B 2009}) Paula the painter had just enough paint for $30$ identically sized rooms. Unfortunately, on the way to work, three cans of paint fell off her truck, so she had only enough paint for $25$ rooms. How many cans of paint did she use for the $25$ rooms?


 $\textbf{(A)}\ 10 \qquad \textbf{(B)}\ 12 \qquad \textbf{(C)}\ 15 \qquad \textbf{(D)}\ 18 \qquad \textbf{(E)}\ 25$

 \textbf{Answer: }C



\textbf{Solution 1}

Losing three cans of paint corresponds to being able to paint five fewer rooms. So $\frac 35 \cdot 25 = \boxed{15}$.  The answer is $\mathrm{(C)}$.

 



\textbf{Solution 2 (proportionality)}

If $p$ is the number of cans,
 \[\frac{p}{30} = \frac{p-3}{25} \implies \frac{p}{p-3} = \frac{6}{5} \implies 5p = 6p-18 \implies p=18\].

 Thus, $p-3 = \boxed{15}$.

 


	\item (\textit{AMC 12 12B 2009}) Twenty percent less than $60$ is one-third more than what number?


 $\textbf{(A)}\ 16\qquad \textbf{(B)}\ 30\qquad \textbf{(C)}\ 32\qquad \textbf{(D)}\ 36\qquad \textbf{(E)}\ 48$

 \textbf{Answer: }D



\textbf{Solution}

Twenty percent less than 60 is $\frac 45 \cdot 60 = 48$. One-third more than a number \textit{n} is $\frac 43n$. Therefore $\frac 43n = 48$ and the number is $\boxed {36}$. The answer is $\mathrm{(D)}$.

 


	\item (\textit{AMC 12 12B 2009}) A rectangular yard contains two flower beds in the shape of congruent isosceles right triangles. The remainder of the yard has a trapezoidal shape, as shown. The parallel sides of the trapezoid have lengths $15$ and $25$ meters. What fraction of the yard is occupied by the flower beds?


 
                    \begin{figure}[H]
                        \centering
                        \includegraphics[width=0.5\textwidth]{Problem/AMC_12_12B_2009_4_Problem_Figure_1.png}
                    \end{figure}
                    
$\textbf{(A)}\ \frac18\qquad \textbf{(B)}\ \frac16\qquad \textbf{(C)}\ \frac15\qquad \textbf{(D)}\ \frac14\qquad \textbf{(E)}\ \frac13$

 \textbf{Answer: }C



\textbf{Solution 1}

Each triangle has leg length $\frac 12 \cdot (25 - 15) = 5$ meters and area $\frac 12 \cdot 5^2 = \frac {25}{2}$ square meters.  Thus the flower beds have a total area of $25$ square meters.  The entire yard has length $25$ m and width $5$ m, so its area is $125$ square meters.  The fraction of the yard occupied by the flower beds is $\frac {25}{125} = \boxed{\frac15}$.  The answer is $\mathrm{(C)}$.

 



\textbf{Solution 2}

The length of each triangle is $\frac{25-15}{2}=5$ meters. By translating then rotating the top right triangle so that it forms a square with the top left triangle, the ratio we desire is literally the ratio of the $5*5=25$ and the total area, which is $5*25=125$. Alternatively we can use side length ratios, but each way we get $\frac{25}{125}=\frac{1}{5}$. Select $\mathrm{(C)}$.

 



\textbf{Solution 3}

Let the trapezoid be defined $ABCD$ in clockwise starting from lower left vertex and the rectangle $AEFD$ exist such that $B$ and $C$ lie on $EF$. Then, drop the altitude from $B$ and $C$ so it intersects at $G$ and $H$. Then, $GH = BC = 15$ and thus $AG = HD = 5$. The triangles are 45-45-90 and thus conjoining both we have a square of side length 5 whose area is 25, and thus the ration is $\frac{25}{25 \cdot 5} = \boxed{\frac{1}{5}}$.

 


	\item (\textit{AMC 12 12B 2009}) Kiana has two older twin brothers. The product of their ages is $128$. What is the sum of their three ages?


 $\textbf{(A)}\ 10\qquad \textbf{(B)}\ 12\qquad \textbf{(C)}\ 16\qquad \textbf{(D)}\ 18\qquad \textbf{(E)}\ 24$

 \textbf{Answer: }D



\textbf{Solution}

The age of each person is a factor of $128 = 2^7$.  So the twins could be $2^0 = 1, 2^1 = 2, 2^2 = 4, 2^3 = 8$ years of age and, consequently Kiana could be $128$, $32$, $8$, or $2$ years old, respectively.  Because Kiana is younger than her brothers, she must be $2$ years old. So the sum of their ages is $2 + 8 + 8 = \boxed{18}$.  The answer is $\mathrm{(D)}$.

 


	\item (\textit{AMC 12 12B 2009}) By inserting parentheses, it is possible to give the expression

 \[2\times3 + 4\times5\]
several values. How many different values can be obtained?


 $\textbf{(A)}\ 2\qquad \textbf{(B)}\ 3\qquad \textbf{(C)}\ 4\qquad \textbf{(D)}\ 5\qquad \textbf{(E)}\ 6$

 \textbf{Answer: }C



\textbf{Solution}

The three operations can be performed on any of $3! = 6$ orders.  However, if the addition is performed either first or last, then multiplying in either order produces the same result. So at most four distinct values can be obtained.  It is easy to check that the values of the four expressions
 \begin{align*} (2\times3) + (4\times5) &= 26,\\ (2\times3 + 4)\times5 &= 50,\\ 2\times(3 + 4\times5) &= 46,\\ 2\times(3 + 4)\times5 &= 70 \end{align*}
are in fact all distinct.  So the answer is $\boxed{4}$, which is choice $\mathrm{(C)}$.

 



\textbf{Solution 2}

We use constructive counting. From $2 \times 3 + 4 \times 5$ there are four ways to place the "(". Then, there are four ways to place the ")". That accounts for 16 plausible cases.

 \textbf{Lemma:} Each value is repeated four times each. We claim the answer is 4.

 \textbf{Proof of the Lemma:}

 Very straightforward. For the single $(2) \times 3 + 4 \times 5$ obviously we see it is repeated four times. For the double, we have that $(2 \times 3) + 4 \times 5 = 2 \times 3 + (4 \times 5)$ is the same. A counting argument easily shows that if one places each of the parentheses such that it follows the Law of PEMDAS, then we obtain cases where we have a repeated sum. We handwave the rest of the proof as an interesting exercise for the reader, but the proof can be built upon PEMDAS.

 \textbf{2009 AMC 10B P7}

 The answer is $\boxed{4}$.

 


	\item (\textit{AMC 12 12B 2009}) In a certain year the price of gasoline rose by $20\%$ during January, fell by $20\%$ during February, rose by $25\%$ during March, and fell by $x\%$ during April. The price of gasoline at the end of April was the same as it had been at the beginning of January. To the nearest integer, what is $x$?


 $\textbf{(A)}\ 12\qquad \textbf{(B)}\ 17\qquad \textbf{(C)}\ 20\qquad \textbf{(D)}\ 25\qquad \textbf{(E)}\ 35$

 \textbf{Answer: }B



\textbf{Solution}

Let $p$ be the price at the beginning of January.  The price at the end of March was $(1.2)(0.8)(1.25)p = 1.2p.$ Because the price at the end of April was $p$, the price decreased by $0.2p$ during April, and the percent decrease was
 \[x = 100 \cdot \frac{0.2p}{1.2p} = \frac {100}{6} \approx 16.7.\]So to the nearest integer $x$ is $\boxed{17}$.  The answer is $\mathrm{(B)}$.

 



\textbf{Solution 2}

Without loss of generality, we can assume the price at the beginning of January was $\$100$. 

 When it rose by $20\%$, it became $\$120$, when it fell by $20\%$, it became $\$96$, and when it rose by $25\%$, it became $\$120$ again. 

 In order for the price at the end of April to be the same as it was at the beginning of January ($\$100$), the price must decrease by $\$20$. 

 20 is $\frac{1}{6}th$ of 120, and $\frac{1}{6} \approx 0.167 \approx 17\%$ So to the nearest integer, $x = 17$ and the answer is $\boxed{\textbf{(B) } 17}$. ~azc1027

 


	\item (\textit{AMC 12 12B 2009}) When a bucket is two-thirds full of water, the bucket and water weigh $a$ kilograms. When the bucket is one-half full of water the total weight is $b$ kilograms. In terms of $a$ and $b$, what is the total weight in kilograms when the bucket is full of water?


 $\textbf{(A)}\ \frac23a + \frac13b\qquad \textbf{(B)}\ \frac32a - \frac12b\qquad \textbf{(C)}\ \frac32a + b \qquad \textbf{(D)}\ \frac32a + 2b\qquad \textbf{(E)}\ 3a - 2b$

 \textbf{Answer: }E



\textbf{Solution 1}

Let $x$ be the weight of the bucket and let $y$ be the weight of the water in a full bucket.  Then we are given that $x + \frac 23y = a$ and $x + \frac 12y = b$.  Hence $\frac 16y = a-b$, so $y = 6a-6b$.  Thus $x = b - \frac 12 (6a-6b) = -3a + 4b$.  Finally $x + y = \boxed {3a-2b}$. The answer is $\mathrm{(E)}$.

 



\textbf{Solution 2}

Imagine that we take three buckets of the first type, to get rid of the fraction. We will have three buckets and two buckets' worth of water.

 On the other hand, if we take two buckets of the second type, we will have two buckets and enough water to fill one bucket.

 The difference between these is exactly one bucket full of water, hence the answer is $3a-2b$.

 



\textbf{Solution 3}

We are looking for an expression of the form $xa + yb$.

 We must have $x+y=1$, as the desired result contains exactly one bucket. Also, we must have $\frac 23 x + \frac 12 y = 1$, as the desired result contains exactly one bucket of water.

 At this moment, it is easiest to check that only the options (A), (B), and (E) satisfy $x+y=1$, and out of these only (E) satisfies the second equation.

 Alternatively, we can directly solve the system, getting $x=3$ and $y=-2$.

 



\textbf{Solution 4}

Since $a$ is one bucket plus two-thirds of the total amount of water, and $b$ is one bucket plus one-half of the total amount of water, $a-b$ would equal $\cancel{\text{bucket}}+\frac23\cdot\text{water}-(\cancel{\text{bucket}}+\frac12\cdot\text{water})=\frac16\cdot\text{water}$. Therefore, $a-b$ is one-sixth of the total mass of the water. Starting from $a$, we add two one-sixths of the total amount of the water to become full. Therefore, the full bucket of water is $a+2(a-b)=a+2a-2b=3a-2b\Rightarrow\textbf{(E)}$

 -\href{https://artofproblemsolving.com/wiki/index.php?title=User:Sweetmango77}{SweetMango77}

 



\textbf{Solution 5 (Using answer choices+Very quick bash)}

WLOG let the bucket hold 6 units of water. Then b=4, a=3. Then test out all the answer choices to arrive at the answer choice (E)~PEKKA

 


	\item (\textit{AMC 12 12B 2009}) Triangle $ABC$ has vertices $A = (3,0)$, $B = (0,3)$, and $C$, where $C$ is on the line $x + y = 7$. What is the area of $\triangle ABC$?


 $\textbf{(A)}\ 6\qquad \textbf{(B)}\ 8\qquad \textbf{(C)}\ 10\qquad \textbf{(D)}\ 12\qquad \textbf{(E)}\ 14$

 \textbf{Answer: }A



\textbf{Solution 1}

Because the line $x + y = 7$ is parallel to $\overline {AB}$, the area of $\triangle ABC$ is independent of the location of $C$ on that line.  Therefore it may be assumed that $C$ is $(7,0)$.  In that case the triangle has base $AC = 4$ and altitude $3$, so its area is $\frac 12 \cdot 4 \cdot 3 = \boxed {6}$.

 



\textbf{Solution 2}

The base of the triangle is $AB = \sqrt{3^2 + 3^2} = 3\sqrt 2$.  Its altitude is the distance between the point $A$ and the parallel line $x + y = 7$, which is $\frac 4{\sqrt 2} = 2\sqrt 2$. Therefore its area is $\frac 12 \cdot 3\sqrt 2 \cdot 2\sqrt 2 = \boxed{6}$.  The answer is $\mathrm{(A)}$.

 
                    \begin{figure}[H]
                        \centering
                        \includegraphics[width=0.5\textwidth]{Solution/AMC_12_12B_2009_9_Solution_Figure_1.png}
                    \end{figure}
                    




\textbf{Solution 3}

By Shoelace, our area is:
 \[\frac {1}{2} \cdot |(3\cdot 3 + 0 \cdot y+ 0 \cdot x)-(0*0+3(x+y))|.\]We know $x+y=7$ so we get:
 \[\frac {1}{2} \cdot |9-21|=\boxed 6\]




\textbf{Solution 4}

WLOG, let the coordinates of $C$ be $(3,4)$ , or any coordinate, for that matter. Applying the shoelace formula, we get the area as $\boxed 6$.

 


	\item (\textit{AMC 12 12B 2009}) A particular $12$-hour digital clock displays the hour and minute of a day. Unfortunately, whenever it is supposed to display a $1$, it mistakenly displays a $9$. For example, when it is 1:16 PM the clock incorrectly shows 9:96 PM. What fraction of the day will the clock show the correct time?


 $\textbf{(A)}\ \frac12\qquad \textbf{(B)}\ \frac58\qquad \textbf{(C)}\ \frac34\qquad \textbf{(D)}\ \frac56\qquad \textbf{(E)}\ \frac {9}{10}$

 \textbf{Answer: }A



\textbf{Solution 1}

The clock will display the incorrect time for the entire hours of $1, 10, 11$ and $12$.  So the correct hour is displayed $\frac 23$ of the time.  The minutes will not display correctly whenever either the tens digit or the ones digit is a $1$, so the minutes that will not display correctly are $10, 11, 12, \dots, 19$ and $01, 21, 31, 41,$ and $51$.  This amounts to fifteen of the sixty possible minutes for any given hour.  Hence the fraction of the day that the clock shows the correct time is $\frac 23 \cdot \left(1 - \frac {15}{60}\right) = \frac 23 \cdot \frac 34  = \boxed{\frac 12}$.  The answer is $\mathrm{(A)}$.

 



\textbf{Solution 2}

The required fraction is the number of correct times divided by the total times. There are 60 minutes in an hour and 12 hours on a clock, so there are 720 total times.

 We count the correct times directly; let a correct time be $x:yz$, where $x$ is a number from 1 to 12 and $y$ and $z$ are digits, where $y<6$. There are 8 values of $x$ that will display the correct time: 2, 3, 4, 5, 6, 7, 8, and 9. There are five values of $y$ that will display the correct time: 0, 2, 3, 4, and 5. There are nine values of $z$ that will display the correct time: 0, 2, 3, 4, 5, 6, 7, 8, and 9. Therefore there are $8\cdot 5\cdot 9=40\cdot 9=360$ correct times.

 Therefore the required fraction is $\frac{360}{720}=\frac{1}{2}\Rightarrow \boxed{\mathrm{(A)}}$.

 


 


	\item (\textit{AMC 12 12B 2009}) On Monday, Millie puts a quart of seeds, $25\%$ of which are millet, into a bird feeder.  On each successive day she adds another quart of the same mix of seeds without removing any seeds that are left.  Each day the birds eat only $25\%$ of the millet in the feeder, but they eat all of the other seeds.  On which day, just after Millie has placed the seeds, will the birds find that more than half the seeds in the feeder are millet?


 $\textbf{(A)}\ \text{Tuesday}\qquad \textbf{(B)}\ \text{Wednesday}\qquad \textbf{(C)}\ \text{Thursday} \qquad \textbf{(D)}\ \text{Friday}\qquad \textbf{(E)}\ \text{Saturday}$

 \textbf{Answer: }D



\textbf{Solution 1}

On Monday, day 1, the birds find $\frac 14$ quart of millet in the feeder.  On Tuesday they find
 \[\frac 14 + \frac 34 \cdot \frac 14\]quarts of millet.  On Wednesday, day 3, they find
 \[\frac 14 + \frac 34 \cdot \frac 14 + \left(\frac34\right)^2 \cdot \frac 14\]quarts of millet.  The number of quarts of millet they find on day $n$ is
 \[\frac 14 + \frac 34 \cdot \frac 14 + \left(\frac34\right)^2 \cdot \frac 14 + \cdots + \left(\frac 34\right)^{n-1} \cdot \frac 14 = \frac {(\frac 14)(1 - (\frac 34)^n)}{1 - \frac 34} = 1 - \left(\frac 34\right)^n .\]The birds always find $\frac 34$ quart of other seeds, so more than half the seeds are millet if $1 - \left(\frac 34\right)^n > \frac 34$, that is, when $\left(\frac 34\right)^n < \frac 14$.  Because $\left(\frac 34\right)^4 = \frac {81}{256} > \frac 14$ and $\left(\frac 34\right)^5 = \frac {243}{1024} < \frac 14$, this will first occur on day $5$ which is $\boxed {\text{Friday}}$. The answer is $\mathrm{(D)}$.

 



\textbf{Solution 2 (Brute Force)}

Notice that, every day, there is always $\frac34$ quarts of other seeds in the feeder, so we simply need to calculate when the quantity of millet seeds is greater than $\frac34$ quarts. Every day that elapses, a fourth of the millet is taken away (eaten by birds), and then a fourth of a quart is added (by Millie). Thus, if $n$ was the quantity of millet seeds the previous day, then on the next day, the quantity is $\dfrac34n+\dfrac14=\dfrac14\left(3n+1\right).$ 
Monday: $\dfrac14\quad$

Tuesday: $\dfrac14\left(3\cdot\dfrac14+1\right)=\dfrac{7}{16}$

Wednesday: $\dfrac14\left(3\cdot\dfrac{7}{16}+1\right)=\dfrac{37}{64}$

Thursday: $\dfrac14\left(3\cdot\dfrac{37}{64}+1\right)=\dfrac{175}{256}$

Friday: $\dfrac14\left(3\cdot\dfrac{175}{256}+1\right)=\dfrac{781}{1024}>\dfrac{768}{1024}=\dfrac34.$

Brute forcing, we see that it takes $4$ days. Thus, the answer is $\boxed{\text{Friday}}$, or $\boxed{\text{(D)}}$. ~ aidan0626, lucaswujc, & Technodoggo
Note: Yes there is a pattern to it, $1-\left(\frac{3}{4}\right)^x$ where $x$ is the number of days, as shown in Solution 1. This solution is only if you didn't notice the pattern. -aidan0626

 



\textbf{Solution 3}

Let $a_n$ denote the the amount of millet that the birds find on day $n$. Here $n\ge1$ and $n=1$ corresponds to Monday. Similarly, let $b_n$ denote the amount of other seeds that the birds find on day $n$. Notice that $a_n$ and $b_n$ satisfy the following:

 \begin{cases}a\_1=\frac{1}{4}\\a\_n=\frac{3}{4}a\_{n-1}+\frac{1}{4}\quad (n>1)\end{cases}

 


 
 \[b_n=\frac{3}{4}\]

 
We can solve for $a_n$ using difference equations to get $a_n=-\left(\frac{3}{4}\right)^n+1$. We want to find the least $n$ for which $a_n>\frac{1}{2}(a_n+b_n)$. In other words, $a_n>b_n$. Substituting in what we know about $a_n$ and $b_n$ we get:

 

 \[-\left(\frac{3}{4}\right)^n+1>\frac{3}{4}\]

 \[\left(\frac{3}{4}\right)^n<\frac{1}{4}\]

 \[3^n<4^{n-1}\]

 \[n\ge5\]

 
Hence, the answer is $\boxed{\text{Friday}}$, or $\boxed{\text{(D)}}$.

 


	\item (\textit{AMC 12 12B 2009}) The fifth and eighth terms of a geometric sequence of real numbers are $7!$ and $8!$ respectively. What is the first term?


 $\textbf{(A)}\ 60\qquad \textbf{(B)}\ 75\qquad \textbf{(C)}\ 120\qquad \textbf{(D)}\ 225\qquad \textbf{(E)}\ 315$

 \textbf{Answer: }E



\textbf{Solution}

Let the $n$th term of the series be $ar^{n-1}$.  Because
 \[\frac {8!}{7!} = \frac {ar^7}{ar^4} = r^3 = 8,\]it follows that $r = 2$ and the first term is $a = \frac {7!}{r^4} = \frac {7!}{16} = \boxed{315}$.  The answer is $\mathrm{(E)}$.

 


	\item (\textit{AMC 12 12B 2009}) Triangle $ABC$ has $AB = 13$ and $AC = 15$, and the altitude to $\overline{BC}$ has length $12$. What is the sum of the two possible values of $BC$?


 $\textbf{(A)}\ 15\qquad \textbf{(B)}\ 16\qquad \textbf{(C)}\ 17\qquad \textbf{(D)}\ 18\qquad \textbf{(E)}\ 19$

 \textbf{Answer: }D



\textbf{Solution}

Let $D$ be the foot of the altitude to $\overline{BC}$.  Then $BD = \sqrt {13^2 - 12^2} = 5$ and $DC = \sqrt {15^2 - 12^2} = 9$.  Thus $BC = BD + BC = 5 + 9 = 14$. Otherwise, assume that the triangle is obtuse at angle $B$ then $BC = DC - BD = 9 -5 = 4$.  The sum of the two possible values is $14 + 4 = \boxed{18}$.  The answer is $\mathrm{(D)}$.

 


	\item (\textit{AMC 12 12B 2009}) Five unit squares are arranged in the coordinate plane as shown, with the lower left corner at the origin. The slanted line, extending from $(c,0)$ to $(3,3)$, divides the entire region into two regions of equal area. What is $c$?

                    \begin{figure}[H]
                        \centering
                        \includegraphics[width=0.5\textwidth]{Problem/AMC_12_12B_2009_14_Problem_Figure_1.png}
                    \end{figure}
                    $\textbf{(A)}\ \frac12\qquad \textbf{(B)}\ \frac35\qquad \textbf{(C)}\ \frac23\qquad \textbf{(D)}\ \frac34\qquad \textbf{(E)}\ \frac45$

 \textbf{Answer: }C



\textbf{Solution 1}

For $c\geq 1.5$ the shaded area is at most $1.5$, which is too little. Hence $c<1.5$, and therefore the point $(2,1)$ is indeed inside the shaded part, as shown in the picture.

 Then the area of the shaded part is one less than the area of the triangle with vertices $(c,0)$, $(3,0)$, and $(3,3)$. The area of the entire triangle is $\frac{3(3-c)}2$, therefore the area of the shaded part is $\frac{7-3c}{2}$.

 The entire figure has area $5$, hence we want the shaded part to have area $\frac 52$. Solving for $c$, we get $c=\boxed{\frac 23}$.  The answer is $\mathrm{(C)}$ ( what a coincidence ;) ).

 



\textbf{Solution 2}

The unit square is of area 1, so the five unit squares have area 5.Therefore the shaded space must occupy 2.5. The missing unit square is of area 1, and if reconstituted the original triangle would be of area 3.5.It can then be inferred:$(3-c) * 3 = 7$.

 $3-c=\frac{7}{3}$, so $3-\frac{7}{3}=c$.

 $3-\frac{7}{3} = \frac{9-7}{3} = \frac{2}{3}$. $\boxed{\mathrm{(C)}}$.

 



\textbf{Solution 3}

The shaded space of the object can become a triangle by adding a unit square to its bottom. This triangle has a base length of $3-c$ and a height of $3$. The area of this triangle region is now (using the formula $A=bh/2$ for a triangle)$(9-3c)/2$. But, remember that we have to subtract the area of the extra unit square from this area to get the area of the shaded region. 

 If you add 3 unit squares to the top of the unshaded portion, the figure is now a right trapezoid. Using the formula for a trapezoid, you can derive that its area is $(9+3c)/2$. After you subtract the area of the 3 extra unit squares, we have derived the area of the unshaded portion.

 Since these areas are both equal, we set them equal to each other and solve for $c$. $(9-3c)/2 -1 = (9+3c)/2 -3$. $c$ is now solved to be $\frac{2}{3}$. $\boxed{\mathrm{(C)}}$.

 



\textbf{Solution 4 (Brute Force)}

We are looking for the area of the shaded region to be $\frac 52$. We start by testing $(A) \frac 12$. The area of the shaded region would be $\frac{(3-\frac 12)(3) }{2}-1=\frac {11}{4}$ when $c=\frac 12$. This does not match our wanted answer.

 We try $(B) \frac 35$ next. The area of the shaded region would be $\frac{(3-\frac 35)(3) }{2}-1=\frac {13}{5}$ when $c=\frac 35$. This also does not match our desired answer.

 We then try $(C) \frac 23$. The area of the shaded region would be $\frac{(3-\frac 23)(3) }{2}-1=\frac {5}{2}$ when $c=\frac 23$. Therefore, our answer is $\boxed{\mathrm{(C) \frac 23}}$.

 



\textbf{Solution 5}

Use the shoelace theorem to get the area of the unshaded portion in terms of a. Divide this area by the whole area which is 5. This fraction should be equal to $1/2$. You can easily solve for the answer from there.-mewto

 


	\item (\textit{AMC 12 12B 2009}) Assume $0 < r < 3$.  Below are five equations for $x$.  Which equation has the largest solution $x$?


 $\textbf{(A)}\ 3(1 + r)^x = 7\qquad \textbf{(B)}\ 3(1 + r/10)^x = 7\qquad \textbf{(C)}\ 3(1 + 2r)^x = 7$
$\textbf{(D)}\ 3(1 + \sqrt {r})^x = 7\qquad \textbf{(E)}\ 3(1 + 1/r)^x = 7$

 \textbf{Answer: }B



\textbf{Solution 1}

\textbf{(B)} Intuitively, $x$ will be largest for that option for which the value in the parentheses is smallest.

 Formally, first note that each of the values in parentheses is larger than $1$, which leads to the conclusion that $x$ will be the largest when the value of the parentheses is smallest. Now, each of the options is of the form $3f(r)^x = 7$. This can be rewritten as $x\log f(r) = \log\frac 73$. As $f(r)>1$, we have $\log f(r)>0$. Thus $x$ is the largest for the option for which $\log f(r)$ is smallest. And as $\log f(r)$ is an increasing function, this is the option for which $f(r)$ is smallest.

 We now get the following easier problem: Given that $0<r<3$, find the smallest value in the set $\{ 1+r, 1+r/10, 1+2r, 1+\sqrt r, 1+1/r\}$.

 Clearly $1+r/10$ is smaller than the first and the third option. 

 We have $r^2 < 10$, dividing both sides by $10r$ we get $r/10 < 1/r$. 

 And finally, $r/100 < 1$, therefore $r^2/100 < r$, and as both sides are positive, we can take the square root and get $r/10 < \sqrt r$.

 Thus the answer is $\boxed{\textbf{(B) } 3\left(1+\frac{r}{10}\right)^x = 7}$.

 



\textbf{Solution 2}

As stated in the first solution, $x$ will be largest when the value in the parenthesis is smallest.

 We can plug in a value for $r$ (for instance, $r=2$) and see which parenthesis has the smallest value. Doing so, we see that B is the answer.

 -Made\_in\_2004

 


	\item (\textit{AMC 12 12B 2009}) Trapezoid $ABCD$ has $AD||BC$, $BD = 1$, $\angle DBA = 23^{\circ}$, and $\angle BDC = 46^{\circ}$.  The ratio $BC: AD$ is $9: 5$.  What is $CD$?


 $\textbf{(A)}\ \frac {7}{9}\qquad \textbf{(B)}\ \frac {4}{5}\qquad \textbf{(C)}\ \frac {13}{15} \qquad \textbf{(D)}\ \frac {8}{9}\qquad \textbf{(E)}\ \frac {14}{15}$

 \textbf{Answer: }B



\textbf{Solution 1}

Extend $\overline {AB}$ and $\overline {DC}$ to meet at $E$.  Then

 
 \begin{align*} \angle BED &= 180^{\circ} - \angle EDB - \angle DBE\\ &= 180^{\circ} - 134^{\circ} -23^{\circ} = 23^{\circ}. \end{align*}


 Thus $\triangle BDE$ is isosceles with $DE = BD$.  Because $\overline {AD} \parallel \overline {BC}$, it follows that the triangles $BCE$ and $ADE$ are similar.  Therefore
 \[\frac 95 = \frac {BC}{AD} = \frac {CD + DE}{DE} = \frac {CD}{BD} + 1 = CD + 1,\]so $CD = \boxed{\frac 45}.$



	\item (\textit{AMC 12 12B 2009}) Each face of a cube is given a single narrow stripe painted from the center of one edge to the center of its opposite edge.  The choice of the edge pairing is made at random and independently for each face.  What is the probability that there is a continuous stripe encircling the cube?


 $\textbf{(A)}\ \frac {1}{8}\qquad \textbf{(B)}\ \frac {3}{16}\qquad \textbf{(C)}\ \frac {1}{4} \qquad \textbf{(D)}\ \frac {3}{8}\qquad \textbf{(E)}\ \frac {1}{2}$

 \textbf{Answer: }B



\textbf{Solution 1}

There are two possible stripe orientations for each of the six faces of the cube, so there are $2^6 = 64$ possible stripe combinations.  There are three pairs of parallel faces so, if there is an encircling stripe, then the pair of faces that does not contribute uniquely determines the stripe orientation for the remaining faces.  In addition, the stripe on each face that does not contribute may be oriented in either of two different ways. So a total of $3 \cdot 2 \cdot 2 = 12$ stripe combinations on the cube result in a continuous stripe around the cube.  The required probability is $\frac {12}{64} = \boxed{\frac {3}{16}}$.

 



\textbf{Solution 2}

Without loss of generality, orient the cube so that the stripe on the top face goes from front to back.  There are two mutually exclusive ways for there to be an encircling stripe: either the front, bottom and back faces are painted to complete an encircling stripe with the top face's stripe or the front, right, back and left faces are painted to form an encircling stripe.  The probability of the first case is $\left(\frac 12\right)^3 = \frac 18$, and the probability of the second case is $\left(\frac 12\right)^4 = \frac {1}{16}$.  The cases are disjoint, so the probabilities sum to $\frac 18 + \frac {1}{16} = \boxed {\frac {3}{16}}$.

 



\textbf{Solution 3}

There are three possible orientations of an encircling stripe (as each encircling would leave out 1 pair of opposite faces, and there are 3 pairs in total).  For any one of these to appear, the stripes on the four faces through which the continuous stripe is to pass must be properly aligned.  The probability of each such stripe alignment is $\left(\frac 12\right)^4 = \frac {1}{16}$.  Since there are three such possibilities and they are disjoint, the total probability is $3 \cdot \frac {1}{16} = \boxed{\frac {3}{16}}$.  The answer is $\mathrm{(B)}$.

 



\textbf{Solution 4}

Consider a vertex on the cube and the three faces that are adjacent to that vertex. If no two stripes on those three faces are aligned, then there is no stripe encircling the cube. The probability that the stripes aren't aligned is $\frac{1}{4}$, since for each alignment of one stripe, there is one and only one way to align the other two stripes out of four total possibilities. therefore there is a probability of $\frac{3}{4}$ that two stripes are aligned.

 Now consider the opposing vertex and the three sides adjacent to it. Given the two connected stripes next to our first vertex, we have two more that must be connected to make a continuous stripe. There is a probability of $\left(\frac{1}{2}\right)^2=\frac{1}{4}$ that they are aligned, so there is a probability of $\frac{3}{4}\cdot \frac{1}{4}=\boxed{\frac {3}{16}}$ that there is a continuous stripe.

 


 



\textbf{Solution 5}

Each face has two possible stripes, so there are $2^6=64$ ways in total.

 We need to choose $2$ opposite faces to leave out. Since the second is fixed once we pick the first face, we only need to calculate $\binom{3}{1}=3$ orientations. Next, the two faces also each have 2 orientations, since they don't matter to the stripe. Therefore there are $3\cdot2\cdot2=12$ ways. The probability is $\frac{12}{64}=\frac{3}{16} \Longrightarrow \boxed{\textbf{(B) } \frac{3}{16}}$.

 


	\item (\textit{AMC 12 12B 2009}) Rachel and Robert run on a circular track. Rachel runs counterclockwise and completes a lap every $90$ seconds, and Robert runs clockwise and completes a lap every $80$ seconds. Both start from the start line at the same time. At some random time between $10$ minutes and $11$ minutes after they begin to run, a photographer standing inside the track takes a picture that shows one-fourth of the track, centered on the starting line. What is the probability that both Rachel and Robert are in the picture?


 $\textbf{(A)}\ \frac {1}{16}\qquad \textbf{(B)}\ \frac18\qquad \textbf{(C)}\ \frac {3}{16} \qquad \textbf{(D)}\ \frac14\qquad \textbf{(E)}\ \frac {5}{16}$

 \textbf{Answer: }C



\textbf{Solution}

After $10$ minutes $(600$ seconds$),$ Rachel will have completed $6$ laps and be $30$ seconds from completing her seventh lap. Because Rachel runs one-fourth of a lap in $22.5$ seconds, she will be in the picture between $18.75$ seconds and $41.25$ seconds of the tenth minute.  After 10 minutes, Robert will have completed $7$ laps and be $40$ seconds from completing his eighth lap.  Because Robert runs one-fourth of a lap in $20$ seconds, he will be in the picture between $30$ seconds and $50$ seconds of the tenth minute.  Hence both Rachel and Robert will be in the picture if it is taken between $30$ seconds and $41.25$ seconds of the tenth minute.  So the probability that both runners are in the picture is $\frac {41.25 - 30} {60} = \boxed{\frac {3}{16}}$.  The answer is $\mathrm{(C)}$.

 


	\item (\textit{AMC 12 12B 2009}) For each positive integer $n$, let $f(n) = n^4 - 360n^2 + 400$. What is the sum of all values of $f(n)$ that are prime numbers?


 $\textbf{(A)}\ 794\qquad \textbf{(B)}\ 796\qquad \textbf{(C)}\ 798\qquad \textbf{(D)}\ 800\qquad \textbf{(E)}\ 802$

 \textbf{Answer: }E



\textbf{Solution 1}

To find the answer it was enough to play around with $f$. One can easily find that $f(1)=41$ is a prime, then $f$ becomes negative for $n$ between $2$ and $18$, and then $f(19)=761$ is again a prime number. And as $41 + 761 = 802$ is already the largest option, the answer must be $\boxed{802}$.

 



\textbf{Solution 2}

We will now show a complete solution, with a proof that no other values are prime. 

 Consider the function $g(x) = x^2 - 360x + 400$, then obviously $f(x) = g(x^2)$.

 The roots of $g$ are:
 \[x_{1,2}  = \frac{ 360 \pm \sqrt{ 360^2 - 4\cdot 400 } }2  = 180 \pm 80 \sqrt 5\]

 We can then write $g(x) = (x - 180 - 80\sqrt 5)(x - 180 + 80\sqrt 5)$, and thus $f(x) = (x^2 - 180 - 80\sqrt 5)(x^2 - 180 + 80\sqrt 5)$.

 We would now like to factor the right hand side further, using the formula $(x^2 - y^2) = (x-y)(x+y)$. To do this, we need to express both constants as squares of some other constants. Luckily, we have a pretty good idea how they look like.

 We are looking for rational $a$ and $b$ such that $(a+b\sqrt 5)^2 = 180 + 80\sqrt 5$. Expanding the left hand side and comparing coefficients, we get $ab=40$ and $a^2 + 5b^2 = 180$. We can easily guess (or compute) the solution $a=10$, $b=4$.

 Hence $180 + 80\sqrt 5 = (10 + 4\sqrt 5)^2$, and we can easily verify that also $180 - 80\sqrt 5 = (10 - 4\sqrt 5)^2$.

 We now know the complete factorization of $f(x)$: 

 
 \[f(x) = (x - 10 - 4\sqrt 5)(x + 10 + 4\sqrt 5)(x - 10 + 4\sqrt 5)(x + 10 - 4\sqrt 5)\]

 As the final step, we can now combine the factors in a different way, in order to get rid of the square roots. 

 We have $(x - 10 - 4\sqrt 5)(x - 10 + 4\sqrt 5) = (x-10)^2 - (4\sqrt 5)^2 = x^2 - 20x + 20$,and $(x + 10 - 4\sqrt 5)(x + 10 + 4\sqrt 5) = x^2 + 20x + 20$.

 Hence we obtain the factorization $f(x) = (x^2 - 20x + 20)(x^2 + 20x + 20)$.

 For $x\geq 20$ both terms are positive and larger than one, hence $f(x)$ is not prime. For $1<x<19$ the second factor is positive and the first one is negative, hence $f(x)$ is not a prime. The remaining cases are $x=1$ and $x=19$. In both cases, $f(x)$ is indeed a prime, and their sum is $f(1) + f(19) = 41 + 761 = \boxed{802}$.

 



\textbf{Solution 3}

Instead of doing the hard work, we can try to guess the factorization. One good approach:

 We can make the observation that $f(x)$ looks similar to $(x^2 + 20)^2$ with the exception of the $x^2$ term. In fact, we have $(x^2 + 20)^2 = x^4 + 40x^2 + 400$. But then we notice that it differs from the desired expression by a square: $f(x) = (x^2 + 20)^2 - 400x^2 = (x^2 + 20)^2 - (20x)^2$.

 Now we can use the formula $(x^2 - y^2) = (x-y)(x+y)$ to obtain the same factorization as in the previous solution, without all the work.

 



\textbf{Solution 4}

After arriving at the factorization $f(x) = (x^2 - 20x + 20)(x^2 + 20x + 20)$, a more mathematical approach would be to realize that the second factor is always positive when $x$ is a positive integer. Therefore, in order for $f(x)$ to be prime, the first factor has to be $1$.

 We can set it equal to 1 and solve for $x$: 

 $x^2-20x+20=1$
$x^2-20x+19=0$
$(x-1)(x-19)=0$
$x=1, x=19$

 Substituting these values into the second factor and adding would give the answer.

 \textbf{A way to find the desired factorization} Essentially since we want to find when the function is prime, we shall attempt to factor the polynomial. Notice we have $(n^2+20)^2-400n^2=n^4-360n^2+400$, and difference of squares gives the desired factorization, and we can continue as in any other solution. Furthermore, the motivation from doing $(n^2+20)^2$ is since it matches the leading term in f(x) and the constant term.

 


	\item (\textit{AMC 12 12B 2009}) A convex polyhedron $Q$ has vertices $V_1,V_2,\ldots,V_n$, and $100$ edges. The polyhedron is cut by planes $P_1,P_2,\ldots,P_n$ in such a way that plane $P_k$ cuts only those edges that meet at vertex $V_k$. In addition, no two planes intersect inside or on $Q$. The cuts produce $n$ pyramids and a new polyhedron $R$. How many edges does $R$ have?


 $\textbf{(A)}\ 200\qquad \textbf{(B)}\ 2n\qquad \textbf{(C)}\ 300\qquad \textbf{(D)}\ 400\qquad \textbf{(E)}\ 4n$

 \textbf{Answer: }C



\textbf{Solution 1}

Each edge of $Q$ is cut by two planes, so $R$ has $200$ vertices. Three edges of $R$ meet at each vertex, so $R$ has a total of $\frac 12 \cdot 3 \cdot 200 = \boxed {300}$ edges.

 



\textbf{Solution 2}

At each vertex, as many new edges are created by this process as there are original edges meeting at that vertex.  Thus the total number of new edges is the total number of endpoints of the original edges, which is $200$.  A middle portion of each original edge is also present in $R$, so $R$ has a total of $100 + 200 = \boxed {300}$ edges.

 



\textbf{Solution 3}

Euler's Polyhedron Formula applied to $Q$ gives $n - 100 + F = 2$, where F is the number of faces of $Q$.  Each edge of $Q$ is cut by two planes, so $R$ has $200$ vertices.  Each cut by a plane $P_k$ creates an additional face on $R$, so Euler's Polyhedron Formula applied to $R$ gives $200 - E + (F+n) = 2$, where $E$ is the number of edges of $R$.  Subtracting the first equation from the second gives $300 - E = 0$, whence $E = \boxed {(C)300}$.

 



\textbf{Solution 4}

Each edge connects two points. The plane cuts that edge so it splits into $2$ at each end (like two legs) for a total of $4$ new edges.

 
                    \begin{figure}[H]
                        \centering
                        \includegraphics[width=0.5\textwidth]{Solution/AMC_12_12B_2009_20_Solution_Figure_1.png}
                    \end{figure}
                    

 But because each new edge is shared by an adjacent original edge cut similarly, the additional edges are overcounted $\times 2$.

 
                    \begin{figure}[H]
                        \centering
                        \includegraphics[width=0.5\textwidth]{Solution/AMC_12_12B_2009_20_Solution_Figure_2.png}
                    \end{figure}
                    

 Since there are $100$ edges to start with, $400/2=200$ new edges result. So there are $100+200=\boxed{\textbf{(C) }300}$ edges in the figure.

 



\textbf{Solution 5}

The question specifies the slices create as many pyramids as there are vertices, implying each vertex owns 4 edge ends. There are twice as many edge-ends as there are edges, and $2 \cdot 100 = 200$.

 $\frac{200}{4} = 50$, so there are $50$ vertices.

 The base of a pyramid has 4 edges, so each sliced vertex would add four edges to $R$.

 $100 + 4 \cdot 50$ = $\boxed{\textbf{(C) } 300}$

 \textbf{Explanation} I doubt anyone has the mental capacity to imagine $Q$ in its full complexity in their mind. So, we begin by examining a simpler structure: a cube.

 
Let $E_C$ denote the number of edges of a polyhedron $C$. Let $V_C$ denote the number of vertices of $C$. Let $B_C$ denote the number of edges that meet at each vertex.

 
Suppose $C$ is a cube. Then

 
$E_C = 12$, $V_C = 8$, and $B_C = 3$

 
(A cube has $12$ edges, $8$ vertices, and $3$ edges meet per vertex)

 
Notice that

 
$V_C*B_C = 2E_C$ ($8*3 = 2*12$)

 
Why? Well, imagine an ant walking along an edge of a cube. When it reaches the end of the edge, it will be standing on a vertex (by definition). But, what if it had walked the other way? Well, it would still be on the same edge, but it would reach the other endpoint of the edge, another vertex. Therefore, we can say that a polyhedron with $E$ edges has $2E$ total endpoints. In the case of $C$, a cube, this means a cube has $2*12 = 24$ total endpoints. But a cube only has $8$ vertices, not $24$... This doesn't make any sense... Wait a minute, we're dealing with a polyhedron, which means multiple edges will be connected to the same vertex. In the case of a cube, $3$ edges meet at every vertex. So we have $24 \div 3$, or $8$ vertices. The math works out. Let's apply this to $Q$, our mystery polyhedron.

 
We are told that $E_Q = 100$, $V_Q = n$ and $B_Q$ isn't given explicitly. So we'll let $B_Q = b$. Now our equation reads

 
$bn = 200$

 
What does this mean, exactly? Well $Q$ has $100$ edges, so it must has $200$ endpoints (endpoints, not vertices). So if we can figure out either $n$, the number of vertices, or $b$, the number of edges per vertex, we can figure out the other variable. Hmmmm... The problem doesn't give us any clues about either. How frustrating. Let's leave this alone and consider the other information given to us.

 
We are also told that "the polyhedron is cut by planes $P_1, P_2,...,P_n$" (at this point I imagine $Q$ is a giant block of cheese, potato, etc.), and also that "plane $P_k$ cuts only those edges that meet at $V_k$". What does this mean?

 
Well, let's go back to our cube. Notice that there are exactly as many planes as vertices. Think of the planes as knives (or knife strokes) cutting into a big cube of cheese. We're cutting the cheese $8$ times, since there are $8$ vertices. Also, whenever we cut into the cheese, we have to direct the cuts near a vertex, cutting into only the edges connected to that vertex. How do we do that? Well, since we can't cut into any other edges, we're gonna have to settle for cutting corners (see what I did there? ...Ok, I admit it was cheesy... see what I did there, too?😂). We have to shave off just the corners of the cheese block, because this way it won't interfere with any of the other edges.

 
The problem also tells us that none of the planes overlap. In our cheese analogy, this means that every time we cut, we have to make sure the cuts are always into new cheese, and not into previously cut areas (all of the cuts are distinct, and don't overlap). Sure enough, this is in agreement with the information that exactly $n$ pyramids are cut off. Let's go back to $Q$.

 
$Q$ has $n$ vertices, so we're gonna make $n$ cuts. The question we want to answer is: "After all of the cuts are made, how many total edges will $R$ (the figure that remains), have?". Well, we're gonna have to figure out either $n$ or $b$... or do we? We're told that $n$ pyramids were made, but it doesn't specify exactly how many sides the pyramid has. So... that implies that it doesn't matter, or else the problem wouldn't be solvable. Let's suppose $b = 4$ ($4$ edges meet at every vertex). Then we would have $200 \div 4 = 50$ vertices, and we would make $50$ cuts. But how many more edges does it yield? (We are told that the planes don't overlap, so the original edges are all present.) Well, if we cut off the corner of a cube, it will create $3$ more edges... which is coincident with the fact that $3$ edges meet at every vertex. So for every vertex, we gain exactly $b$ new edges. Rewriting the equation to find the number of vertices, we have

 
$v = \frac{200}{b}$

 
But we just said that for every vertex/cut, we gain $b$ more edges. So the total number of new edges is $v*b \Rightarrow \frac{200}{b}*b \Rightarrow 200$. It doesn't matter how many vertices $Q$ has, or how many edges meet at every vertex. Given the conditions in the problem, we will always gain $200$ more edges, regardless of the values of $n$, or $b$. We have $200$ new edges $+$ $100$ old edges to make $300$ edges $\Rightarrow \boxed{\text{C}}$.

 


	\item (\textit{AMC 12 12B 2009}) Ten women sit in $10$ seats in a line. All of the $10$ get up and then reseat themselves using all $10$ seats, each sitting in the seat she was in before or a seat next to the one she occupied before. In how many ways can the women be reseated?


 $\textbf{(A)}\ 89\qquad \textbf{(B)}\ 90\qquad \textbf{(C)}\ 120\qquad \textbf{(D)}\ 2^{10}\qquad \textbf{(E)}\ 2^2 3^8$

 \textbf{Answer: }A



\textbf{Solution 1}

Notice that either a woman stays in her own seat after the rearrangement, or two adjacent women swap places. Thus, our answer is counting the number of ways to arrange 1x1 and 2x1 blocks to form a 1x10 rectangle. This can be done via casework depending on the number of 2x1 blocks. The cases of 0, 1, 2, 3, 4, 5 2x1 blocks correspond to 10, 8, 6, 4, 2, 0 1x1 blocks, and so the sum of the cases is
 \[\binom{10}{0} + \binom{9}{1} + \binom{8}{2} + \binom{7}{3} + \binom{6}{4} + \binom{5}{5} = 1 + 9 + 28 + 35 + 15 + 1 = \boxed{89}.\]




\textbf{Solution 1 Shortcut}

Recall that the number of ways to arrange 1x1 and 2x1 blocks to form a 1xn rectangle results in \href{https://artofproblemsolving.com/wiki/index.php?title=Fibonacci\_numbers}{Fibonacci numbers}. Clearly, $\boxed{89}$ is the only fibonacci number, so no calculation is needed for this problem.

 



\textbf{Solution 2}

Let $S_n$ be the number of possible seating arrangements with $n$ women. Consider $n \ge 3,$ and focus on the rightmost woman. If she returns back to her seat, then there are $S_{n-1}$ ways to seat the remaining $n-1$ women. If she sits in the second to last seat, then the woman who previously sat there must now sit at the rightmost seat. This gives us $S_{n-2}$ ways to seat the other $n-2$ women, so we obtain the recursion
 \[S_n = S_{n-1}+S_{n-2}.\]

 Starting with $S_1=1$ and $S_2=2,$ we can calculate $S_{10}=\boxed{89}.$

 \textbf{Clarification of Solution 2} The seating possibilities of woman \#$10$ become the two cases which we work out. $S_n$ was defined to be the number of different seating arrangements with $n$ women.

 In the first "case" if woman \#$10$ sits in the seat \#$10$, this leads to a similar scenario, but with $9$ women instead. That means that for this case, there are a total of $S_{9}$ possible arrangements. We don't know how many exactly, but we are able to define it in terms of $S$.

 During the second "case", woman \#$10$ sits in seat \#$9$. This time, woman \#9 must go to seat \#$10$, as she is the only other person who can go there. This leaves us with $8$ women, and we again represent this in terms of $S \Rightarrow S_8$.

 Therefore, we can write $S_{10}$ in terms of $S_8$ and $S_9$, like so:

 
 \[S_{10} = S_8 + S_9.\]

 We can then generalize this to say

 
 \[S_n = S_{n-1}+S_{n-2}.\]

 Calculating $S_1 = 1$ and $S_2 =2,$ then following the recursive rule from above, we get $S_{10} = 89 \Rightarrow \boxed{\text{A}}$.

 


	\item (\textit{AMC 12 12B 2009}) Parallelogram $ABCD$ has area $1,\!000,\!000$. Vertex $A$ is at $(0,0)$ and all other vertices are in the first quadrant. Vertices $B$ and $D$ are lattice points on the lines $y = x$ and $y = kx$ for some integer $k > 1$, respectively. How many such parallelograms are there?


 $\textbf{(A)}\ 49\qquad \textbf{(B)}\ 720\qquad \textbf{(C)}\ 784\qquad \textbf{(D)}\ 2009\qquad \textbf{(E)}\ 2048$

 \textbf{Answer: }C



\textbf{Solution 1}

The area of any parallelogram $ABCD$ can be computed as the size of the vector product of $\overrightarrow{AB}$ and $\overrightarrow{AD}$.

 In our setting where $A=(0,0)$, $B=(s,s)$, and $D=(t,kt)$ this is simply $s\cdot kt - s\cdot t = (k-1)st$. 

 In other words, we need to count the triples of integers $(k,s,t)$ where $k>1$, $s,t>0$ and $(k-1)st = 1,\!000,\!000 = 2^6 5^6$.

 These can be counted as follows: We have $6$ identical red balls (representing powers of $2$), $6$ blue balls (representing powers of $5$), and three labeled urns (representing the factors $k-1$, $s$, and $t$). The red balls can be distributed in ${8\choose 2} = 28$ ways, and for each of these ways, the blue balls can then also be distributed in $28$ ways. (See \href{https://artofproblemsolving.com/wiki/index.php?title=Distinguishability}{Distinguishability} for a more detailed explanation.)

 Thus there are exactly $28^2 = 784$ ways how to break $1,\!000,\!000$ into three positive integer factors, and for each of them we get a single parallelogram. Hence the number of valid parallelograms is $784 \longrightarrow \boxed{C}$.

 



\textbf{Solution 2}

Without the vector product the area of $ABCD$ can be computed for example as follows: If $B=(s,s)$ and $D=(t,kt)$, then clearly $C=(s+t,s+kt)$. Let $B'=(s,0)$, $C'=(s+t,0)$ and $D'=(t,0)$ be the orthogonal projections of $B$, $C$, and $D$ onto the $x$ axis. Let $[P]$ denote the area of the polygon $P$. We can then compute:

 
 \begin{align*} [ABCD] & = [ADD'] + [DD'C'C] - [BB'C'C] - [ABB'] \\ & = \frac{kt^2}2 + \frac{s(s+2kt)}2 - \frac{t(2s+kt)}2 - \frac{s^2}2 \\ & = kst - st \\ & = (k-1)st. \end{align*}
The remainder of the solution is the same as the above.

 


 



\textbf{Solution 3}

We know that $A$ is $(0, 0)$. Since $B$ is on the line $y=x$, let it be represented by the point $(b, b)$. Similarly, let $D$ be $(d, kd)$. Since this is a parallelogram, sides $\overline{AD}$ and $\overline{BC}$ are parallel. Therefore, the distance and relative position of $D$ to $A$ is equivalent to that of $C$ to $B$ (if we take the translation of $A$ to $D$ and apply it to $B$, we will get the coordinates of $C$). This yields $C (b+d, b+kd)$. Using the \href{https://artofproblemsolving.com/wiki/index.php?title=Shoelace\_Theorem}{Shoelace Theorem} we get

 
$1,000,000 = \frac{1}{2}|\left(b(b+kd) + (b+d)(kd)\right) - \left(b(b+d) + (b+kd)d\right)|$
$\Rightarrow 2,000,000 = |2kbd - 2bd|$
$\Rightarrow 1,000,000 = |kbd - bd|$

 
Since $k > 1, kbd > bd$. The equation becomes

 
$1,000,000 = (k-1)bd$
$\Rightarrow \frac{1,000,000}{k-1} = bd$

 
Since $k$ must be a positive integer greater than $1$, we know $k-1$ will be a positive integer. We also know that $bd$ is an integer, so $k-1$ must be a factor of $1,000,000$. Therefore $bd$ will also be a factor of $1,000,000$.

 Notice that $1,000,000 = 10^6 = 2^6*5^6$.

 Let $b$ be $2^x5^y$ such that $x, y$ are integers on the interval $[0, 6]$.

 Let $d$ be $2^w5^z$ such that $w, z$ are integers, $x+w \le 6$, and $y+z \le 6$.

 For a pair $(x, y)$, there are $7-x$ possibilities for $w$ and $7-y$ possibilites for $z$ ($d$ doesn't have to be the co-factor of $1,000,000$, it just can't be big enough such that $bd > 1,000,000$), for a total of $(7-x)(7-y)$ possibilities. So we want

 
$\sum_{k=0, i=0}^6 (7-k)(7-i)$
$\left(\text{the sum of the number of possible pairs}(w, z) \text{ for all pairs}(k , i)\text{ for } k[0, 6]\text{ and } i[0, 6]\right)$

 
Notice that if we "fix" the value of $k$, at, say $6$, then run through all of the values of $i$, change the value of $k$ to $5$, and run through all of the values of $i$ again, and so on until we exhaust all $49$ combinations of $(k, i)$ we get something like this:

 
$1*1 + 1*2 + ... + 1*6 + 1*7 + 2*1 + 2*2 + ... + 2*6 + 2*7 + ..... + 7*1 + 7*2 + ... + 7*6 + 7*7$

 
which can be rewritten

 
$1(1+2+...+7)+2(1+2+...+7)+.....+7(1+2+...+7)$
$\Rightarrow (1+2+...+7)(1+2+...+7)$
$\Rightarrow 28^2$
$\Rightarrow 784$

 
So there are $784$ possible sets of coordinates $B,$ $C$, and $D \Rightarrow \boxed{\text{C}}$.

 Note: this solution could be greatly simplified by using the Shoelace Formula on the triangle $ABD$, which we know has half the area of the parallelogram. This eliminates the need to find the coordinates of point $C$. 

 (Notational note: I'm not sure if the notation for double index summation is correct or even applicable in the context of this problem. If someone could fix the notation so that it is correct, or replace it without changing the general content of this solution, that would be great. If the notation is correct, then just delete this footnote)

 


	\item (\textit{AMC 12 12B 2009}) A region $S$ in the complex plane is defined by

 \[S = \{x + iy: - 1\le x\le1, - 1\le y\le1\}.\]
A complex number $z = x + iy$ is chosen uniformly at random from $S$. What is the probability that $\left(\frac34 + \frac34i\right)z$ is also in $S$?


 $\textbf{(A)}\ \frac12\qquad \textbf{(B)}\ \frac23\qquad \textbf{(C)}\ \frac34\qquad \textbf{(D)}\ \frac79\qquad \textbf{(E)}\ \frac78$

 \textbf{Answer: }D



\textbf{Solution 1}

First, turn $\frac34 + \frac34i$ into polar form as $\frac{3\sqrt{2}}{4}e^{\frac{\pi}{4}i}$. Restated using geometric probabilities, we are trying to find the portion of a square enlarged by a factor of $\frac{3\sqrt{2}}{4}$ and rotated $45$ degrees that lies within the original square. This skips all the absolute values required before. Finish with the symmetry method stated above.

 -asdf334

 


 



\textbf{Solution 2}

We multiply $z$ and $\left(\frac{3}{4}+\frac{3}{4}i\right)$ to get 
 \[\left(\frac{3}{4}x-\frac{3}{4}y\right)+\left(\frac{3}{4}xi+\frac{3}{4}yi\right).\] Since we want to find the probability that this number is in $S$, we need the real and complex coefficients of this number to be less than or equal to $1$ or greater than or equal to $-1.$ This gives us the equations 
 \[-1\le \frac{3}{4}x-\frac{3}{4}y \le 1\] and 
 \[-1\le \frac{3}{4}x+\frac{3}{4}y\le 1.\] Now, we see that we can solve this by graphing. We can graph our barriers $-1\le x\le 1$ and $-1\le y\le 1$ to form a $2$ by $2$ square centered at the origin. Graphing our two equations gives us the four lines 
 \[x-y=\frac{4}{3},\] 
 \[x-y=-\frac{4}{3},\] 
 \[x+y=\frac{4}{3},\] 
 \[x+y=-\frac{4}{3}.\] The square that is formed is the region that satisfies these four equations. Now, the barriers and this square gives us an octagon as the desired region. The area of this octagon is the total area of the square minus the 4 small triangles on each corner, each with $\frac{2}{9}$ area. Therefore, the octagon has area of $\frac{28}{9}.$ Finally, to find the probability of it working, we find the area of the octagon divided by the area of the entire square which is $\frac{\frac{28}{9}}{4}=\frac{7}{9}$ or $\boxed{D}.$

 -jeteagle

 


	\item (\textit{AMC 12 12B 2009}) For how many values of $x$ in $[0,\pi]$ is $\sin^{ - 1}(\sin 6x) = \cos^{ - 1}(\cos x)$?
Note: The functions $\sin^{ - 1} = \arcsin$ and $\cos^{ - 1} = \arccos$ denote inverse trigonometric functions.


 $\textbf{(A)}\ 3\qquad \textbf{(B)}\ 4\qquad \textbf{(C)}\ 5\qquad \textbf{(D)}\ 6\qquad \textbf{(E)}\ 7$

 \textbf{Answer: }B



\textbf{Solution 1}

First of all, we have to agree on the range of $\sin^{-1}$ and $\cos^{-1}$. This should have been a part of the problem statement -- but as it is missing, we will assume the most common definition: $\forall x: -\pi/2 \leq \sin^{-1}(x) \leq \pi/2$ and $\forall x: 0\leq \cos^{-1}(x) \leq \pi$.

 Hence we get that $\forall x\in[0,\pi]: \cos^{ - 1}(\cos x) = x$, thus our equation simplifies to $\sin^{ - 1}(\sin 6x) = x$.

 Consider the function $f(x) = \sin^{ - 1}(\sin 6x)  - x$. We are looking for roots of $f$ on $[0,\pi]$.

 By analyzing properties of $\sin$ and $\sin^{-1}$ (or by computing the derivative of $f$) one can discover the following properties of $f$:

 
            \begin{itemize}
        \item $f(0)=0$.
        \item $f$ is increasing and then decreasing on $[0,\pi/6]$.
        \item $f$ is decreasing and then increasing on $[\pi/6,2\pi/6]$.
        \item $f$ is increasing and then decreasing on $[2\pi/6,3\pi/6]$.
            \end{itemize}
            For $x=\pi/6$ we have $f(x) = \sin^{-1} (\sin \pi) - \pi/6 = -\pi/6 < 0$. Hence $f$ has exactly one root on $(0,\pi/6)$.

 For $x=2\pi/6$ we have $f(x) = \sin^{-1} (\sin 2\pi) - 2\pi/6 = -2\pi/6 < 0$. Hence $f$ is negative on the entire interval $[\pi/6,2\pi/6]$.

 Now note that $\forall t: \sin^{-1}(t) \leq \pi/2$. Hence for $x > 3\pi/6$ we have $f(x) < 0$, and we can easily check that $f(3\pi/6)<0$ as well.

 Thus the only unknown part of $f$ is the interval $(2\pi/6,3\pi/6)$. On this interval, $f$ is negative in both endpoints, and we know that it is first increasing and then decreasing. Hence there can be zero, one, or two roots on this interval.

 To prove that there are two roots, it is enough to find any $x$ from this interval such that $f(x)>0$.

 A good guess is its midpoint, $x=5\pi/12$, where the function $\sin^{-1}(\sin 6x)$ has its local maximum. We can evaluate:$f(x) = \sin^{-1}(\sin 5\pi/2) - 5\pi/12 = \pi/2 - 5\pi/12 = \pi/12 > 0$.

 Summary: The function $f$ has $\boxed{\textbf{(B) }4}$ roots on $[0,\pi]$: the first one is $0$, the second one is in $(0,\pi/6)$, and the last two are in $(2\pi/6,3\pi/6)$.

 Actual solutions are $x=0$, $x=\pi/7$, $x=2\pi/5$, and $x=3\pi/7$.

 



\textbf{Solution 2}

Since $0\leq \cos^{-1}(a) \leq \pi$ for all $a$, the equation reduces to $\sin^{-1}(\sin(6x)) = x$. Since $-\pi/2 \leq \sin^{-1}(a) \leq \pi/2$ for all $a$, $0 \leq x \leq \pi/2$. To make the problem easier, we will measure angles in degrees. We will consider each sixth of the interval $[0, 90]$.

 For $0 \leq x \leq 15$, $6x$ is in the first quadrant. Thus, $\sin^{-1}(\sin(6x)) = 6x$. Setting this equal to $x$ yields the solution $x = 0$.

 For $15 \leq x \leq 30$, $6x$ is in the second quadrant. Thus, $\sin^{-1}(\sin(6x)) = 180 - 6x$. This yields the solution $x = \frac{180}7$.

 For $30 \leq x \leq 45$, $6x$ is in the third quadrant. Thus, $\sin^{-1}(\sin(6x)) = 180 - 6x$. As $\frac{180}{7}$ is not on the interval $30 \leq x \leq 45$, this yields no solution.

 For $45 \leq x \leq 60$, $6x$ is in the fourth quadrant. Thus, $\sin^{-1}(\sin(6x)) = 6x - 360$. As $72$ is not on the interval $45 \leq x \leq 60$, this yields no solution.

 For $60 \leq x \leq 75$, $6x$ is in the first quadrant plus a full revolution. Thus, $\sin^{-1}(\sin(6x)) = 6x - 360$. This yields the solution $x = 72$.

 For $75 \leq x \leq 90$, $6x$ is in the second quadrant plus a full revolution. Thus $\sin^{-1}(\sin(6x)) = 540 - 6x$. This yields the solution $x = \frac{540}7$.

 There are $\boxed{\textbf{(B) }4}$ solutions, $x=0$, $x=\pi/7$, $x=2\pi/5$, and $x=3\pi/7$.

 



\textbf{Solution 3}

Algebraically, the inverse function of a function should just cancel out, leaving $6x=x$. However, upon inspection we find that the graphs of these "inverse function of the function" equations are also periodic, like their normal trig function counterparts, due to the fact that inverse trig functions will never return any angle value higher than $2\pi$. But instead of a smooth wave, these graphs are made up of zigzags with slope $1$ and $-1$. Trying a few values, we see that

 
$y = \arcsin{\sin{x}}$

 
peaks at $\frac{\pi}{2}$ and

 
$y = \arccos{\cos{x}}$

 
peaks at $\pi$

 
But we want the graph of $y = \arcsin{(\sin6{x})}$, which has a period of $\frac{\pi}{3}$ instead of $2\pi$. So this means the interval $[0, \pi]$ will show $3$ periods instead of $\frac{1}{2}$ of a period. Visually it would be lines with slopes $6$ and $-6$. Using the graph paper given to us, we plot out the two equations according to the above and we see that they intersect $4$ times $\Rightarrow \boxed{\text{B}}$.

 



\textbf{Solution 4}

The most conventional domain and range for $\cos^{-1} x$ is $x \in [-1, 1]$ and $\cos^{-1} x \in [0, \pi]$. For $\cos x$, $x \in [0, \pi]$ and $\cos x \in [-1, 1]$. Since the domain of $\cos x$ is equal to the range of $\cos^{-1} x$, and the range of $\cos x$ is equal to the domain of $\cos^{-1} x$, $\cos^{-1} (\cos x) = x$.

 Meaning that $\sin^{ - 1}(\sin 6x) = x$, $\sin x = \sin (6x)$.

 By graphing $\sin x$ and $\sin 6x$ together, it can be seen that there are 7 intersections for $x \in [0, \pi]$.

 
                    \begin{figure}[H]
                        \centering
                        \includegraphics[width=0.5\textwidth]{Solution/AMC_12_12B_2009_24_Solution_Figure_1.png}
                    \end{figure}
                    However, the range of $\sin^{-1} x$ is $[-\frac{\pi}{2}, \frac{\pi}{2}]$. Therefore $x \in [0, \frac{\pi}{2}]$ for $x = \sin^{ - 1}(\sin 6x)$. There are $\boxed{\textbf{(B) }4}$ solutions for $x \in [0, \frac{\pi}{2}]$.

 The $4$ solutions are respectively:

 $x = 6x$, $x = 0$
$x + 2 \pi = 6x$, $x = \frac25 \pi$
$\pi - x = 6x$, $x = \frac{\pi}{7}$
$3 \pi - x = 6x$, $x = \frac37 \pi$




\textbf{Solution 5 (Intervals)}

Since the range of arccos(x) is just [0,$\pi$], arccos(cos(x)) = x Now, we have arcsin(sin(6x)) = x. The range of arcsin is [-$\pi$/2,$\pi$/2], but x is restricted to [0,$\pi$]. This means that x is restricted to  the intersection of both intervals or [0,$\pi$/2]. On the interval for x of [0,$\pi$/3), 6x is in the interval [0,2$\pi$). This means that sin(x) can only be equal to sin(6x) at a maximum of 2 values of x. Either x = 6x, or x = $\pi$-6x. For the first equation, x = 0; for the second equation, x = $\pi$/7. Both are on the interval [0,$\pi$/3), so we have exactly 2 solutions for x on the interval [0,$\pi$/2]. For the interval for x of [$\pi$/3,$\pi$/2), 6x is on the interval [2$\pi$,3$\pi$). This also indicates that sin(x) can only be equal to sin(6x) at a maximum of 2 values of x. Either x = 6x-2$\pi$, or x = $\pi$ - (6x-2$\pi$). For the first equation, x = 2$\pi$/5; for the second equation, x = 3$\pi$/7. Both are on the interval [$\pi$/3, $\pi$/2], so we have exactly 2 solutions for x on this interval as well. Finally, we see that at x = $\pi$/2 sin(6x) = 0 and sin(x) = 1, so x = $\pi$/2 is not a solution. Therefore, there are 4 solutions to sin(6x) = sin(x), namely 0, $\pi$/7, 2$\pi$/5, and 3$\pi$/7. The Answer is (B).

 


	\item (\textit{AMC 12 12B 2009}) The set $G$ is defined by the points $(x,y)$ with integer coordinates, $3\le|x|\le7$, $3\le|y|\le7$. How many squares of side length at least $6$ have their four vertices in $G$?

                    \begin{figure}[H]
                        \centering
                        \includegraphics[width=0.5\textwidth]{Problem/AMC_12_12B_2009_25_Problem_Figure_1.png}
                    \end{figure}
                    $\textbf{(A)}\ 125\qquad \textbf{(B)}\ 150\qquad \textbf{(C)}\ 175\qquad \textbf{(D)}\ 200\qquad \textbf{(E)}\ 225$

 \textbf{Answer: }E



\textbf{Solution}

We need to find a reasonably easy way to count the squares. 

 First, obviously the maximum distance between two points in the same quadrant is $4\sqrt 2 < 6$, hence each square has exactly one vertex in each quadrant.

 Given any square, we can circumscribe another axes-parallel square around it. In the picture below, the original square is red and the circumscribed one is blue.

 
                    \begin{figure}[H]
                        \centering
                        \includegraphics[width=0.5\textwidth]{Solution/AMC_12_12B_2009_25_Solution_Figure_1.png}
                    \end{figure}
                    

 Let's now consider the opposite direction. Assume that we picked the blue square, how many different red squares do share it?

 Answering this question is not as simple as it may seem. Consider the picture below. It shows all three red squares that share the same blue square. In addition, the picture shows a green square that is not valid, as two of its vertices are in bad locations.

 
                    \begin{figure}[H]
                        \centering
                        \includegraphics[width=0.5\textwidth]{Solution/AMC_12_12B_2009_25_Solution_Figure_2.png}
                    \end{figure}
                    

 The size of the blue square can range from $6\times 6$ to $14\times 14$, and for the intermediate sizes, there is more than one valid placement. We will now examine the cases one after another. Also, we can use symmetry to reduce the number of cases.

 \begin{verbatim}
size  upper\_right  solutions  symmetries  total    6        (3,3)          1           1      1    7        (3,3)          1           4      4    8        (3,3)          1           4      4    8        (3,4)          1           4      4    8        (4,4)          3           1      3    9        (3,3)          1           4      4    9        (3,4)          1           8      8    9        (4,4)          3           4     12   10        (3,3)          1           4      4   10        (3,4)          1           8      8   10        (3,5)          1           4      4   10        (4,4)          3           4     12   10        (4,5)          3           4     12   10        (5,5)          5           1      5   11        (4,4)          3           4     12   11        (4,5)          3           8     24   11        (5,5)          5           4     20   12        (5,5)          5           4     20   12        (5,6)          5           4     20   12        (6,6)          7           1      7   13        (6,6)          7           4     28   14        (7,7)          9           1      9
\end{verbatim}

Summing the last column, we get that the answer is $\boxed{225}$.

 



\textbf{Solution 2}

This is based on a clever bijection given in \href{https://artofproblemsolving.comhttp://web2.slc.qc.ca/sh/Contest/AMC12\_2009B-S.pdf}{this page}.

 Consider any square $ABCD$ where all four vertices are in $G$, and the side length is at least $6$, so the four vertices must lie in distinct quadrants (Same proof as in solution 1). Without loss of generality, assume that $A,B,C,D$ are in the first, second, third, fourth quadrant. Then we consider the following mapping:

 
 \[A \to A' = A\]
 \[B \to B' = B + (10,0)\]
 \[C \to C' = C + (10,10)\]
 \[D \to D' = D + (0,10)\]

 Then the new points $A'$, $B'$, $C'$, $D'$ are either being the same point or forming a square in $G_1 = G \cap \{ x>0, y>0 \}$, a 5x5 grid. 

 Conversely, for any point in $G_1$, it can be reversed to a square $ABCD$; however, for any square in $G_1$, there are four possible squares $ABCD$ that were mapped to them. Therefore the number of possible squares $ABCD$ is equal to $25 + 4N$, where $N$ is the number of squares inscribed in $G_1$.

 Moreover by the same idea in solution 1, each square (with sides parallel or slanted to the axes) in a $G_1$ can be inscribed in a square in $G_1$, with sides parallel to one of the axes, call it "standard square". Noticing that each standard square of side length $a$ corresponds to $a$ inscribed squares,  and that there are $(5-a)^2$ number of standard squares of side length $a$, we have

 
 \[N = \sum_{a=1}^{4} a(5-a)^2 = 1\cdot 16 + 2\cdot 9 + 3 \cdot 4 + 4 \cdot 1 = 16+18+12+4=50\]

 So the answer is $25+4\cdot 50 = 225$

 Motivation(by williamgolly, solution by MAA):Homotheties

 



\textbf{Solution 3(If you have no time)}

Start with Solution 2: the number of possible squares $ABCD$ is equal to $25 + 4N$. Notice there are only 2 options ($A$, and $E$) which have integer solutions for N. Option $A$ corresponds to $N=25$, and option $E$ corresponds to $N=50$. At this point, guessing between the 2 choices is already worth it, but there are already 16 unit squares (solutions to N), meaning that the answer is more probably $E$.

 Edit:There are 125 squares you can make w/out tilted ones, so the answer is more than 125, hence $E$.

 



\textbf{Solution 4/Explanation}

Consider any square that meets the requirements described in the problem. Then, take the vertices of the square and translate them to the first quadrant (This is the "mapping" described in Solution 2). For example, consider a square with vertices $(7, 7), (-7, 7), (-7, -7),$ and $(7, -7)$:

 

                    \begin{figure}[H]
                        \centering
                        \includegraphics[width=0.5\textwidth]{Solution/AMC_12_12B_2009_25_Solution_Figure_3.png}
                    \end{figure}
                    

 
After following the mapping described in Solution 2, the square looks like this:

 

                    \begin{figure}[H]
                        \centering
                        \includegraphics[width=0.5\textwidth]{Solution/AMC_12_12B_2009_25_Solution_Figure_4.png}
                    \end{figure}
                    

 
The position of each vertex within their corresponding grid has not changed.For example, the point $(-7, 7)$ is still the top-left point in a grid, albeit a change in quadrant. Trying this out with a couple of other squares, we see that the following property holds:

 
 \[\text{"For any and all squares with vertices in each of the four quadrants, the 'mapped' version is also a square."}\]

 
Therefore the logical inverse is true:

 

 \[\text{"For any 'mapped' square with all four vertices in the first quadrant,}\]
 \[\text{there exists at least one corresponding 'unmapped' square with a vertex in each of the four quadrants."}\]

 
But how many "unmapped" squares, to be exact? 

 
This might seem complicated at first, but with some intuitive thinking, we realize that there are exactly $4$ "unmapped" squares that correspond with a "mapped" square. This is because given a "mapped" square, there are $4$ choices for the vertex that will remain in the first quadrant; but once that point is chosen, there is only $1$ distribution of the other $3$ vertices that will result in a square. So, we want four times the number of squares we can make in the first quadrant grid.

 
We divide our counting method into two cases: squares with side length $0$ after mapping (which means all four vertices are in the same position relative to their own grids) and squares with side length $1-4$ after mapping.

 
Case 1: There are $25$ such squares of length $0$ (this is equivalent to counting the number of points on the grid). However, in this scenario, all of the vertices have been mapped onto the same point. So instead of $4$ choices for the first quadrant vertex, there is only one. Subsequently there are  only $25$ such squares that correspond to them.

 
Case 2: Let a square with sides parallel to the axes be known as $A$ squares. These $A$ squares can have side length $1, 2, 3,$ or $4$. However, the number of $A$ squares possible depends on the side length. For example, there is only $1$ possible $A$ square of side length $4$, but $16$ squares of side length $1$. To be exact, there are $(5-s)^2$ possible $A$ squares of side length $s$. So, the total number of $A$ squares is

 
$\sum_{s=1}^4 (5-s)^2$

 
But what about "tilted" squares? Notice that "tilted squares" can always be inscribed (drawn within) another, bigger square. Let a square inscribed within an $A$ square be called a $B$ square. How many $B$ squares are there? Well, this also depends on the side length. We only want squares whose vertices are lattice points (integer value coordinates), so the number of $B$ squares should increase along with side length. We defined $B$ squares to be inscribed within $A$ squares, so we can say that all $B$ squares have their vertices on the side on an $A$ square. Consider an $A$ square with side length $4$. There are $3$ other lattice points along the side of the $A$ square, not counting the vertices. Therefore, we can say that there are $s-1$ possible $B$ squares for every $A$ square with side length $s$. We can multiply $(s-1)$ times the number of $A$ squares to get the number of $B$ squares. This is 

 
$\sum_{s=1}^4 (s-1)(5-s)^2$

 
total $B$ squares. But we need to add these two quantities to get the number of squares for Case 2:

 
$\sum_{s=1}^4 (s-1)(5-s)^2 + \sum_{s=1}^4 (5-s)^2$

 
By distributive property, the expression becomes

 
$\sum_{s=1}^4 s(5-s)^2$

 
Solving, we get $50$ "mapped" squares, both $A$ and $B$. Multiplying this by $4$ to get the corresponding number of "unmapped" squares, then adding to get the number of squares for Case 1, we get

 
$50*4 + 25 = 225 \Rightarrow \boxed{\text{E}}$.

 \textbf{Lemma (For Solution 2)} We may prove, using vectors, that two regular $n$-sided polygons create a a new regular $n$-sided polygon on a rectangular coordinate system.

 Let the two polygons be $M_1 M_2 ... M_n$ and $N_1 N_2 ... N_n$.

 Let $\omega = e^{i\frac{2\pi}{n}}$ and $a = \vec{M_1 M_2}$.We have the respective vector expressions of each point of $M$, considering $M_1$ as a starting point.

 $\begin{cases} 	\vec{M_1} = \vec{M_1} \\ 	\vec{M_2} = \vec{M_1}+\vec{a} \\ 	\vec{M_3} = \vec{M_1}+\vec{a}+\omega\vec{a} \\         \vec{M_4} = \vec{M_1}+\vec{a}+\omega\vec{a}+\vec{a}+\omega^2 \vec{a} \\         ...... \\         \vec{M_n} = \vec{M_1}+(1+\omega+\omega^2+......+\omega^{n-1})\vec{a} \end{cases}$

 The case is similar with $N$.

 Adding up each corresponding point of $M$ and $N$, we get the points for the new polygon.

 $\begin{cases} 	\vec{M_1} + \vec{N_1} = \vec{M_1} + \vec{N_1} \\ 	\vec{M_2} + \vec{N_2} = \vec{M_1} + \vec{N_1} + (\vec{a} + \vec{b})\\         \vec{M_3} + \vec{N_3} = \vec{M_1} + \vec{N_1} + (1 + \omega)(\vec{a} + \vec{b})         \vec{M_4} + \vec{N_4} = \vec{M_1} + \vec{N_1} + (1 + \omega + \omega^2)(\vec{a} + \vec{b})\\         ...... \\         \vec{M_n} + \vec{N_n} = \vec{M_1} + \vec{N_1} + (1+\omega+\omega^2+......+\omega^{n-1})(\vec{a} + \vec{b})  \end{cases}$

 See how the expression of the new polygon's points is similar to $M$ and $N$? We conclude that, adding up a $n$-sided regular polygon with an "initial point" on $\vec{M_1}$ and a side-length of $\left| \vec{a}\right|$ with another $n$-sided regular polygon with an "initial point" on $\vec{N_1}$ and a side-length of $\left| \vec{b}\right|$ creates a new $n$-sided regular polygon with its "initial point" on $\vec{M_1}+\vec{N_1}$ and a side-length of $\left| \vec{a}+\vec{b}\right|$.

 Solution $2$ is a special case, where the sum of vertices of $2$ squares is a square on a complex plane.

 


\end{enumerate}